\section{The auxiliary dihedral representation at every coefficient level}

Let $K\subset\Omega$ be the imaginary quadratic field in the original
seed problem, and let $\ell$ be an odd prime prime to its discriminant.
Assume that every decomposition group over $\ell$ fixes $K$. Write
$G=\operatorname{Gal}(\Omega/\mathbb Q)$, $H=\operatorname{Gal}(\Omega/K)$,
$E=\overline{\mathbb Q}_\ell$, and $\varepsilon$ for the full cyclotomic
character, also on restriction to $H$. Fix the comparison
$\iota:E\simeq\mathbb C$ and the complex conjugation $c$ from the
hypotheses. We give the construction underlying $\mathrm{TauExAlt}$.
The one field drawn from $\mathrm{Af.Hypotheses}'$ by the formal
construction is its global Artin reciprocity law. CIP1 already
proves this law from classical class field theory, with the required
convention that a uniformizer maps to inverse arithmetic Frobenius.

% BEGIN D66T1
\paragraph{D66T1. A Hecke character with the prescribed finite restriction.}
Put $N=|\operatorname{disc} K|\ge3$. The quadratic character
$\chi:(\mathbb Z/N)^\times\to\{\pm1\}$ of $K$ is odd.
Indeed, for the fundamental discriminant $D=\operatorname{disc} K$,
the quadratic Gauss sum for the Kronecker symbol $(D/\cdot)$ has
square $D$ and is multiplied by $(D/u)$ under $\zeta_N\mapsto\zeta_N^u$.
Thus $K=\mathbb Q(\sqrt D)$ affords that character, and complex
conjugation acts nontrivially. Fix one embedding $e:K\to\mathbb C$.
There is a finite character of $(\mathcal O_K/N)^\times$ agreeing
with $\chi$ on rational units modulo $N$ and with $e$ on global units.

To see compatibility, all units $u$ of an imaginary quadratic field
are roots of unity: their conjugate absolute values are both one and
the ring of integers is a lattice in $\mathbb C$. If $u$ is congruent
to an integer modulo $N$, then $u-\bar u$ belongs to $N\mathcal O_K$.
The two conjugates of $(u-\bar u)/N$ have absolute value at most
$2/N<1$. Its integral norm therefore vanishes and $u=\pm1$.
Rational residues inject since $N\mathcal O_K\cap\mathbb Z=N\mathbb Z$.
The prescriptions agree on both signs because $\chi(-1)=-1$.
They define a character on the generated subgroup, which extends to
the finite abelian group by choosing an appropriate root in
$\mathbb C^\times$ each time a generator is adjoined. Its values on
the prescribed subgroup have order prime to $\ell$: an $\ell$th root
of unity in $K$ would make $\ell$ ramify in $K$, by total ramification
of $\mathbb Q_\ell(\zeta_\ell)/\mathbb Q_\ell$. Project the values of
the extension to their prime-to-$\ell$ parts. The two prescriptions
are preserved. Denote this finite character by $b$ and its exponent by $m_0$.

On ideles with all finite coordinates units prescribe
$z_\infty^{-1}b(z_f\bmod N)$, in the coordinate given by $e$.
Reduction uses the local factors of $N\mathcal O_K$ and the Chinese
remainder theorem. On principal ideles prescribe one. The common
domain is the diagonal global units, on which the two factors cancel.
These formulas define a homomorphism on the product of the two
subgroups. It extends to all ideles: their quotient by that product
is the ideal class group, and adjoining representatives successively
only requires roots in $\mathbb C^\times$. The resulting character
$Y$ is continuous since it agrees on an open subgroup with the
displayed continuous formula. Its infinity exponents are $(-1,0)$
at $(e,\bar e)$. It is unramified at $w\nmid N$, and all finite local
units have image killed by $m_0$. At a rational prime $q\nmid N$
with $(q)=w$ prime in $K$, divide the principal idele $q$ by the
idele equal to $q$ at $w$ and one elsewhere. The quotient has finite
unit coordinates, reduces to $q\bmod N$, and has infinite coordinate
$q$. Hence
\[
                       Y_w(q)=q\,\chi(q)^{-1}.
\]
% END D66T1

% BEGIN D66T2
\paragraph{D66T2. The Galois character and its values on squares.}
Take the Artin map of CIP1. Classical global reciprocity identifies
the idele class group modulo its identity component with the
abelianized Galois group in the given closure $\Omega$. Inverting
arithmetic reciprocity gives an open quotient onto $G_K^{\rm ab}$
with kernel the closure generated by principal ideles and the
connected archimedean subgroup. Local reciprocity carries units
onto the image of inertia and uniformizers to inverse Frobenius
modulo inertia; the closure of each local image is the image of
its decomposition group. These are precisely the five original
conditions checked in CIP1. For a Hecke character of
algebraic type $a$, divide out its archimedean monomial, transport
the resulting smooth character by $\iota^{-1}$, and multiply by
the $\ell$-adic monomial
\[
             \prod_{w\mid\ell}\prod_{j:K_w\hookrightarrow E}
                 j(x_w)^{a(\iota j|_K)} .
\]
Here $j$ ranges over continuous embeddings. The smooth character
has open kernel: near one its archimedean restriction is trivial,
and on the profinite finite unit group its continuous complex
character has finite image (a small neighborhood of one in
$\mathbb C^\times$ contains no nontrivial subgroup).
The transferred character is thus continuous even though
$\iota$ is only a field isomorphism. The product is one on connected
archimedean ideles. On a principal idele, the transported inverse
archimedean monomial cancels the product over the $\ell$-adic
embeddings, since each global embedding occurs once. Hence it
descends through Artin to a continuous character with the original
local evaluation law. Apply this to the $Y$ of D66T1 and call the
result $\kappa_0$. The unit part of reciprocity shows
$\kappa_0^{m_0}=1$ on inertia away from $\ell$.

For an inert prime $q\nmid N\ell$ and an arithmetic Frobenius lift
$t$ above $q$, $t^2$ is a Frobenius above the prime $(q)$ of residue
degree two. The inverse-uniformizer convention and D66T1 give
$\kappa_0(t^2)=\chi(q)/q=-\varepsilon(t)^{-1}$. These Frobenius
elements are dense in the coset $G\setminus H$: in any finite Galois
quotient also recording $K$, Chebotarev supplies a prime outside
$N\ell$ in the indicated conjugacy class, and changing its prolongation
selects the specified representative. Both sides of the square
identity are continuous there, so
\[
 \kappa_0(t^2)=-\varepsilon(t)^{-1}\quad(t\notin H),\qquad
 \kappa_0(h)\kappa_0(chc^{-1})=\varepsilon(h)^{-1}\quad(h\in H).
\]
For the second identity use $t=hc$, $c^2=1$ and
$\varepsilon(c)=-1$. Every element outside $H$ is $c$ times an
element of $H$, and a character is invariant under inner conjugation
by $H$. Thus this conjugation identity also holds with any such
element in place of $c$.
% END D66T2

% BEGIN D66T3
\paragraph{D66T3. Split inertia and the change of algebraic type.}
Fix a prolongation $A$ at $\ell$ and its place $w$ in $K$.
The image of the rational decomposition group in
$\operatorname{Gal}(K/\mathbb Q)$ is trivial, so $K_w=\mathbb Q_\ell$.
There is exactly one continuous embedding $j:K_w\to E$.
The two embeddings of $K$ in $\mathbb C$ are $e,\bar e$.
Consequently $a(\iota j|_K)$ is either $-1$ or zero. Every inertia
class in $G_K^{\rm ab}$ at this place is represented by a local unit
$x$. Since $w\nmid N$, the local finite character is trivial,
and the formula of D66T2 gives $j(x)^{a(\iota j|_K)}$.
The same computation with the idele norm, whose type is $1$ at
every embedding, gives $j(x)$. Its avatar is the cyclotomic
character: at each prime away from $\ell$ its arithmetic Frobenius
value is the residue cardinality, since the norm of a uniformizer
is its reciprocal. Chebotarev and continuity identify the
characters globally. Thus on the entire inertia group at $A$,
$\kappa_0=\varepsilon^{-1}$ or $\kappa_0=1$. This proves
$\mathrm{TwistIne}$ for every original $\mathrm{IsTwist}$ witness,
since only its type and unramifiedness at $w$ entered the computation.

More generally, for any number field and any original Artin map,
if $Y$ with type $a$ has avatar $\psi$, the character
$Y^r |\cdot|_{\mathbb A}^{s}$ has type $r a+s$ and avatar
$\psi^r\varepsilon^s$ for all $r,s\ge0$. The archimedean
monomial, the finite local factors and the $\ell$-adic monomial
all multiply term by term. The norm identification just proved
uses the same Frobenius argument over this field, and its
archimedean factor at complex places is the product over both
embeddings. This proves the original $\mathrm{AlgType}$, including
the zero powers, under the hypothesis of $\mathrm{algType\_alt}$.
% END D66T3

% BEGIN D66T4
\paragraph{D66T4. A single character and exponent for all powers ($\mathrm{KapExAlt}$).}
Let $\lambda=\kappa_0^2\varepsilon|_H$. Its conjugate under each
element outside $H$ equals $\lambda^{-1}$ by D66T2, and
$\lambda(t^2)=(-\varepsilon(t)^{-1})^2\varepsilon(t)^2=1$.
The subgroup of $E^\times$ given by $|x-1|_\ell<1$ is open, by
the ultrametric inequality and $|x|_\ell=1$ on it. Its inverse
image under $\lambda$ has finite index in the compact group $H$.
Choose $d>0$ annihilating the finite quotient and set
\[
                      m=d m_0,\qquad \kappa=\lambda^m .
\]
Fix any $j\ge1$, and set $M=mj$. The continuous character
$\kappa^j$ is congruent to one everywhere, is inverted by every
off-field conjugation, and is one on every off-field square.
At every inertia group away from $\ell$ its value is one because
$m_0\mid M$ and $\varepsilon$ is unramified. On each inertia group
above $\ell$, D66T3 gives either $\varepsilon^{M}$ or
$\varepsilon^{-M}$ for the whole group. Finally D66T3 with
$(r,s)=(2M,M)$ supplies an algebraic Hecke character on the
same Artin map; the indices $2M(-1)+M=-M$ and $2M(0)+M=M$
give the two ordered exponents. These are all the fields of
$\mathrm{KapDat}(\iota,K,M,\kappa^j)$. Both $m$ and $\kappa$
were chosen before $j$.
% END D66T4

% BEGIN D66T5
\paragraph{D66T5. One coefficient ring for the entire construction.}
The continuous character $\kappa:H\to E^\times$ has image in
one finite extension of $\mathbb Q_\ell$. Recall the coefficient
argument of RV4. The finite subextensions of $E$ form a countable
family: Krasner approximation of a separable primitive-element
polynomial by rational polynomials gives a root arbitrarily close
to that generator and of the same degree, hence generating the
same embedded field. There are countably many such polynomials and
finitely many roots each. The inverse images of the multiplicative
groups of these fields are closed in $H$ and cover it. By Baire's
theorem one has interior and is an open subgroup. Compactness
gives finitely many cosets; adjoining their character values gives
one finite field $E_1$ containing the full image. The norm of this
compact image is bounded. If a value had norm greater than one
its positive powers would be unbounded, and applying the same
argument to the inverse excludes norm less than one. The integer
ring $O_1$ of $E_1$ inside $E$ is therefore a source coefficient
ring containing all values and their inverses. For every permitted
coefficient ring $O_2\supset O_1$ and each $j\ge1$, all values
of $\kappa^j$ already belong to $O_2^\times$. Their differences
from one lie in its maximal ideal by D66T4 and the valuation
criterion for the integer ring of a finite local field.
% END D66T5

% BEGIN D66T6
\paragraph{D66T6. The integral representation at the specified mark.}
For the general $\mathrm{TauOf}$ construction let $\eta:H\to E^\times$
and $M$ satisfy $\mathrm{KapDat}$ and let its values lie in the
given coefficient ring $O_2$. They are units since the values on
the inverse elements also belong to $O_2$. Fix the original
$g_0\notin H$ and put $J=\operatorname{diag}(1,-1)$. For $h\in H$
write $k=\eta(h)$ and let
\[
 a_h=(k+k^{-1})/2,\qquad b_h=(k-k^{-1})/2,\qquad
 B(h)=\begin{pmatrix}a_h&b_h\\b_h&a_h\end{pmatrix}.
\]
These belong to $O_2$ since $\ell$ is odd. Define $\tau(h)=B(h)$
and $\tau(hg_0)=B(h)J$. In the basis $(1,1)^t,(1,-1)^t$,
$B(h)$ has eigenvalues $k,k^{-1}$, while $J$ interchanges the
two vectors. Hence $B$ is multiplicative and
$B(g_0 h g_0^{-1})=JB(h)J$. The square condition gives
$B(g_0^2)=1=J^2$. These are all the products to check: bringing
one occurrence of $g_0$ past an element of $H$ uses the conjugation
relation, and two occurrences multiply to $g_0^2\in H$.
This proves multiplicativity on both cosets and gives exactly
$\tau(g_0)=J$, even when $g_0$ has order greater than two.

We have $a_h^2-b_h^2=1$, so both $B(h)$ and $J$ preserve the
form $\operatorname{diag}(1,-1)$. On $H$ the diagonal entries are
equal and the off-diagonal entries are equal. On $Hg_0$ the rows
are $(a_h,-b_h),(b_h,-a_h)$, giving the required opposite entries.
The congruence $k\equiv1$ implies $a_h\equiv1$, $b_h\equiv0$.
The two cosets are open, and each entry on each coset is built from
continuous functions of $h$ by operations in $O_2$. The valuation
topology on this finite local integer ring is its maximal-ideal
adic topology, which proves the source entrywise continuity.
On inertia away from $\ell$ inside $H$, $k=1$ and $B(h)=1$.
Finally $a_h+b_h=\eta(h)$ gives the original character recovery
identity on all of $H$.
% END D66T6

% BEGIN D66T7
\paragraph{D66T7. Transport of the integral local signs.}
Use the character and weight $(\eta,M)$ from D66T6.
For every valuation $A$ of $\Omega$ define $s(A)=1$ if
$\eta=\varepsilon^M$ on the subgroup of $H$ lying in rational
inertia at $A$, and $s(A)=-1$ otherwise. At $A$ over $\ell$ the
splitting assumption puts its whole rational decomposition group in
$H$. The two alternatives of $\mathrm{KapDat}$ are mutually
exclusive there. Indeed, inertia in $G_{\mathbb Q_\ell}$ surjects
onto $\mathbb Z_\ell^\times$ through $\varepsilon$: each
$\mathbb Q_\ell(\zeta_{\ell^n})$ is totally ramified
(the shifted cyclotomic polynomial is Eisenstein), its Galois
action realizes all units modulo $\ell^n$, and compactness gives
surjectivity in the inverse limit. The rational decomposition
group at the given prolongation identifies with this local Galois
group and preserves its inertia. Since $M>0$ and $1+\ell$ has
infinite order, $\varepsilon^{2M}$ is nontrivial on this inertia.

If $gA=A'$, conjugation by $g$ carries inertia at $A$ to inertia
at $A'$ and preserves $\varepsilon$. It preserves $\eta$ for
$g\in H$ and inverts it otherwise. Uniqueness of the local sign
therefore gives $s(A')=s(A)$ in the first case and $s(A')=-s(A)$
in the second. For $h$ in rational inertia over $\ell$ the
required integral identity is
\[
                  \tau(h)_{00}+s(A)\tau(h)_{01}
                         =\varepsilon(h)^M .
\]
For $s(A)=1$ the left side is $\eta(h)$, and for $s(A)=-1$ it is
$\eta(h)^{-1}$. The local alternative gives the displayed equality
in both cases. All the signs are defined at every valuation, as
required even where there is no transport or inertia clause.
Together with D66T6, the positivity, continuity of $\eta$ and
the algebraic witness already in $\mathrm{KapDat}$, this proves
the full $\mathrm{TauOf}$ implication.
% END D66T7

% BEGIN D66T8
\paragraph{D66T8. The dihedral supplier for the integral seed ($\mathrm{TauExAlt}$).}
Fix any field, comparison map, prime, quadratic field and
complex conjugation allowed in the hypotheses, and any specified
$g_0\in G\setminus H$. Take the single $m>0$ and character
$\kappa$ of D66T4 and the coefficient ring $O_1$ of D66T5.
For every source coefficient ring $O_2\supset O_1$ and every
positive integer $j$, apply D66T6 and D66T7 to
$(\eta,M)=(\kappa^j,mj)$. The resulting integral representation
and local signs satisfy
\[
              \mathrm{TauDat}(\iota,K,g_0,mj,O_2,\tau,s,\kappa^j).
\]
The eigencharacter on the quadratic field is exactly $\kappa^j$,
the algebraic exponents are $(-mj,mj)$ at the original conjugate
embeddings, and the diagonal normalization is at the specified
element $g_0$ for every enlargement and every power. Thus the first arithmetic supplier in S65A1 works for the
given seed field on the class field theory already reconstructed
in CIP1. The same witnesses prove the source implication with
its stated $\mathrm{Af.Hypotheses}'$ premise.
% END D66T8
